\subsection*{Defining I-Tuples}

An integer tuple, written \co@ituple@, groups several integer values in a fixed order. Its initializer uses braces:

\begin{lstlisting}[language=croqtile,caption={An i-tuple literal}]
ituple a = {1, 2, 3};
b = {4, 5, 6};  // utilize the type inference
\end{lstlisting}

Since Croqtile can infer types from the \emph{initialization expression}, programmers can often omit the \texttt{ituple} keyword. However, without explicit type annotation, an \emph{ituple} variable definition might look similar to an \emph{mdspan} definition if you are not yet familiar with Croqtile. To distinguish them:

\begin{itemize}[leftmargin=1.6em]
\item The \emph{initialization expression} of an \emph{ituple} follows an assignment operation \texttt{=} (like integers), whereas an \emph{mdspan} is initialized after \texttt{:};
\item The \emph{initialization expression} of an \emph{ituple} is enclosed by \texttt{\{\}}, not \texttt{[]} as for \emph{mdspan}.
\end{itemize}

Similar to \texttt{mdspan}, you may enforce rank check for \emph{ituple}s at compile time:

\begin{lstlisting}[language=croqtile,caption={An inconsistent rank}]
ituple<3> a = {1, 2};  // error: inconsistent rank
\end{lstlisting}
